\section*{Chapter \ref{ch:rs:anatomy-of-a-predictor} --- Anatomy of a Predictor}

\begin{solution}{exo:rs:anatomy-of-a-predictor:1}
$1/\sqrt{799} = 0.0354$. A mean of 0.01 reaches $t = 2$ when $0.01\sqrt n/0.0354 = 2$, $n = 50$ independent dates, if the only noise
were the cross-sectional sampling; common shocks make real IC series noisier, and several times as many dates are needed in practice.
\end{solution}

\begin{solution}{exo:rs:anatomy-of-a-predictor:2}
Still 0.02 against the market-excess return: subtracting one number from every return leaves the ranks unchanged. Against
beta-adjusted returns it changes whenever betas differ, because each stock subtracts a different amount; how much depends on how
the \omterm{def:rs:anatomy-of-a-predictor:predictor}{predictor} correlates with beta.
\end{solution}

\begin{solution}{exo:rs:anatomy-of-a-predictor:3}
$0.012/0.08 \times \sqrt{252} = 2.38$.
\end{solution}

\begin{solution}{exo:rs:anatomy-of-a-predictor:4}
The residual keeps $\rho = 0.6$ of the variance: $0.03/\sqrt{0.6} = 0.039$.
\end{solution}

\begin{solution}{exo:rs:anatomy-of-a-predictor:5}
Up to $\sqrt{20} = 4.47$; a naive 6.0 may be as little as $6.0/4.47 = 1.34$.
\end{solution}

\begin{solution}{exo:rs:anatomy-of-a-predictor:6}
Rank IC: robust to outliers in the \omterm{def:rs:anatomy-of-a-predictor:predictor}{predictor} and in returns, so one extreme name does not decide a date. Pearson IC: a portfolio whose
weights are proportional to the \omterm{def:rs:anatomy-of-a-predictor:predictor}{predictor} earns in proportion to the Pearson correlation, magnitudes included, so it measures what a
linear book would make.
\end{solution}

\begin{solution}{exo:rs:anatomy-of-a-predictor:7}
On the tenth of days with the largest absolute market return the rank IC is 0.037 ($t = 5.0$); on the others 0.039 ($t = 17.1$). The
reversal is carried evenly: the simulator reverses a fixed share of each day's specific shock, whatever the market does. On real data
this is the kind of dependence to test, not to assume.
\end{solution}

\begin{solution}{exo:rs:anatomy-of-a-predictor:8}
The 60-day targets of consecutive days overlap by 59 days, so the 1\,260 ICs are nowhere near independent: there are about 21
non-overlapping periods. The proposition bounds the overstatement at $\sqrt{60}$: the honest $t$ may be as low as $14/7.75 = 1.8$. Use
Newey--West with 59 lags, or every 60th date.
\end{solution}

\begin{solution}{pb:rs:anatomy-of-a-predictor:1}
\textbf{1.} 0.0385, $t = 17.8$. \textbf{2.} 0.0074, $t = 2.4$. \textbf{3.} 0.0162 (one-day against five) and 0.0187 (five-day against
one). \textbf{4.} One-day against one day: the simulator reverses a share of each day's specific shock the next day, and each other
pair adds four days of unrelated returns to the \omterm{def:rs:anatomy-of-a-predictor:predictor}{predictor} or to the target. \textbf{5.} 0.0008: nothing. \textbf{6.} Ranks do not change
when one number is subtracted from every return. \textbf{7.} 10\% (the residual keeps 90\%). \textbf{8.} 52\%; $1/\sqrt{0.476} = 1.45$.
\textbf{9.} 0.0043 ($t = 0.8$), 0.0067 ($t = 3.5$), 0.0086 ($t = 3.5$). \textbf{10.} The \omterm{def:rs:anatomy-of-a-predictor:predictor}{predictor}'s own industry component carries no
information about the future in this market, so removing it removes noise from the \omterm{def:rs:anatomy-of-a-predictor:predictor}{predictor} as the target's residual removes noise
from the target. \textbf{11.} 3.85 and 2.39. \textbf{12.} The bound $\sqrt h$ holds for a persistent \omterm{def:rs:anatomy-of-a-predictor:predictor}{predictor}; reversal changes daily, so
the ICs of consecutive dates overlap less. \textbf{13.} $t = 1.31$: valid, but from a fifth of the dates. \textbf{14.} $1/\sqrt{999} =
0.032$. \textbf{15.} 0.040, 6.4, 20.1. \textbf{16.} Under a day: the IC is 0.0385 with day $t+1$ and 0.0008 with day $t+2$.
\textbf{17.} The failure modes and then the IC itself: more liquidity providers compete the reversal away (crowding, chapter~28).
\textbf{18.} \emph{Named result:} for five-day reversal scored daily against overlapping five-day targets the naive $t$ is 3.85 and the
HAC $t$ 2.39; in the strong-industry market, neutralising the target and the \omterm{def:rs:anatomy-of-a-predictor:predictor}{predictor} raises the IC from 0.0043 to 0.0086 and its $t$
from 0.8 to 3.5. \textbf{19.} The \omterm{def:rs:anatomy-of-a-predictor:predictor}{predictor}'s formula and window, the target and its horizon, the normalisation and \omterm{def:rs:anatomy-of-a-predictor:neutralise}{neutralisation},
the universe and period, the IC type, and the standard error used. \textbf{20.} A formula with a target, a horizon, a normalisation and
a \omterm{def:rs:anatomy-of-a-predictor:neutralise}{neutralisation}, of which the formula is the least informative part.
\end{solution}

\begin{solution}{iq:rs:anatomy-of-a-predictor:1}
The cross-sectional correlation between a \omterm{def:rs:anatomy-of-a-predictor:predictor}{predictor} and the return it predicts, computed per date and averaged. For stock returns a
mean IC of 0.02 to 0.05 is good if it is stable; the information ratio of the IC series and its persistence matter as much as its
mean.
\iqlookfor{a definition per date, realistic magnitudes, and attention to stability.}
\end{solution}

\begin{solution}{iq:rs:anatomy-of-a-predictor:2}
Not as stated: consecutive 20-day targets overlap by 19 days, so the naive $t$ can overstate by up to $\sqrt{20} = 4.5$; $t = 8$ may be
about 1.8. Recompute with Newey--West (19 lags) or every 20th date.
\iqlookfor{the overlap correction with a number.}
\end{solution}

\begin{solution}{iq:rs:anatomy-of-a-predictor:3}
To remove a component that carries no information but adds noise (industry moves), and to avoid a book that is really an industry
bet. Not when the \omterm{def:rs:anatomy-of-a-predictor:predictor}{predictor}'s information is at the industry level (industry momentum), or when the book is meant to take industry
risk.
\iqlookfor{\omterm{def:rs:anatomy-of-a-predictor:neutralise}{neutralisation} as a statement about where the information is.}
\end{solution}

\begin{solution}{iq:rs:anatomy-of-a-predictor:4}
Most of its apparent power comes from its correlation with the factors that were removed (industry, beta, size): it is partly a
factor bet. The 0.02 is its stock-specific skill; whether the rest is worth having depends on whether that factor exposure is
rewarded and wanted.
\iqlookfor{the decomposition into factor and specific parts.}
\end{solution}

\begin{solution}{iq:rs:anatomy-of-a-predictor:5}
Pearson, if positions will be proportional to the forecast (the magnitudes matter), checked alongside the rank IC for robustness; rank
if the forecast is used only to order names or if its tails are unreliable.
\iqlookfor{the link between the metric and how the forecast is used.}
\end{solution}

\begin{solution}{iq:rs:anatomy-of-a-predictor:6}
Under no predictability and independent names, the sample correlation of $N$ pairs has variance about $1/(N - 1)$. The mean of $n$
independent dates has standard error $1/\sqrt{n(N - 1)}$; $t = 2$ for a mean of 0.02 with $N = 1\,000$ needs $n = (2/(0.02 \times
31.6))^2 = 10$ dates in theory, but common shocks make the IC's true dispersion several times larger, and in practice hundreds of
dates are needed.
\iqlookfor{the $1/\sqrt{N-1}$ result and awareness that it understates real noise.}
\end{solution}
